\subsection{Hardware Design}
\label{sec:des-hardware}
\subsubsection{Plotter designs and mechanics}
The chosen method of text plotting is a Cartesian gantry system with 2 independent axes of movement with a pen up/down mechanicsm. Complex kinematics equations and error dependencies are avoiding by using 2 isolated x and y movement methods for reproduction \cite{teikari2016}. Linear GT2 timing belts and pulleys are used to convert rotary motion from stepper motors and axles to the desired cartesian movements. Timing belts are used for their reliability, cost effectiveness, and higher speed travel \cite{iqs2024}.

A Cartesian gantry consists of a stationary rail that controls movement in one axis, and a movable rail that controls the other axis of movement. For this system and future explanations, the stationary rails control y-axis movement, while the x-axis rail moves along the y-axis. The following diagram explains axes of the gantry:

\begin{figure}[!htbp]
    \centering
    \includegraphics[width=0.72\linewidth]{Figures/gantry-axis-diagram.png}
    \caption{Cartesian gantry axes and pen up/down mechanism.}
    \label{fig:gantry-axis-diagram}\label{fig:des-gantry-axes}
\end{figure}

The x-axis rail needs to move with stability and consistency, which is why it is supported on both ends by dual-rails on the y-axis. A singular stepper motor was chosen to drive both pulleys on both y-axis rails to ensure synchronicity and save costs. 2 independent stepper motors driving one axis creates synchronicity issues because any difference in timing, missed steps or load torques could offset the movement between the two motors. The challenge with using one motor to drive two timing pulleys, is finding a method to extend the shaft of the motor to connect to both timing pulleys without overloading the motor and applying too much tortion or twist to the new axle. A rigid coupler was used to extend the axle of the stepper motor to a chromed steel rod.

To reduce strain on motors, and allow for increased accuracy and precision, linear bearing were fitted around the chromed rod rails to allow all parts to slide with very little friction.
To verify the stepper motor could handle the weight of the x-axis rail further, the torque required to move the rail at the desired speed was calculated. It is described in Section 3.3.2 with the descriptions of the stepper motor type chosen and its product specifications.


\begin{figure}[!htbp]
    \centering
    \begin{minipage}[c]{0.48\linewidth}
        \centering
        \includegraphics[width=\linewidth,height=0.16\textheight,keepaspectratio]{Figures/cad-component-01.png}
    \end{minipage}\hfill
    \begin{minipage}[c]{0.48\linewidth}
        \centering
        \includegraphics[width=\linewidth,height=0.16\textheight,keepaspectratio]{Figures/cad-component-02.png}
    \end{minipage}\\[0.5em]
    \begin{minipage}[c]{0.48\linewidth}
        \centering
        \includegraphics[width=\linewidth,height=0.16\textheight,keepaspectratio]{Figures/cad-component-03.png}
    \end{minipage}\hfill
    \begin{minipage}[c]{0.48\linewidth}
        \centering
        \includegraphics[width=\linewidth,height=0.16\textheight,keepaspectratio]{Figures/cad-component-04.png}
    \end{minipage}\\[0.5em]
    \begin{minipage}[c]{0.48\linewidth}
        \centering
        \includegraphics[width=\linewidth,height=0.16\textheight,keepaspectratio]{Figures/cad-component-05.png}
    \end{minipage}\hfill
    \begin{minipage}[c]{0.48\linewidth}
        \centering
        \includegraphics[width=\linewidth,height=0.16\textheight,keepaspectratio]{Figures/cad-component-06.png}
    \end{minipage}\\[0.5em]
    \begin{minipage}[c]{0.48\linewidth}
        \centering
        \includegraphics[width=\linewidth,height=0.16\textheight,keepaspectratio]{Figures/cad-component-07.png}
    \end{minipage}\hfill
    \begin{minipage}[c]{0.48\linewidth}
        \centering
        \includegraphics[width=\linewidth,height=0.16\textheight,keepaspectratio]{Figures/cad-component-08.png}
    \end{minipage}
    \caption{CAD models of the gantry components.}
    \label{fig:gantry-cad-components}
\end{figure}


\label{sec:des-camera}
The camera support holds the USB camera over the restrained page so that the image-processing stage receives a consistent view of the writing area. The IMX179 autofocus webcam and page restraint connect to the hardware arrangement shown in the system block diagram.

\subsubsection{Stepper motor and driver}
\label{sec:des-stepper}
A stepper motor does not continuously turn with a PWM input, but rather turns at specific fixed angles each time there is a switch in the windings (the coils of the stepper motor). The A4988 regulates the winding current through a process called chopping, which switches the current supply on the winding until the current rises to a certain point, and then switches off at the target current.  This occurs at high frequencies and effectively keeps the current going to the coils within a specific range so as to not overshoot, even when the motor is not actively stepping. A step is created by slighty adjusting the current value between 2 coils just enough to allow the rotors teeth to make a small jump or 'step'. The stepper motor is perfect for the precise movements required for handwriting reproduction with the gantry, as it produces steps at exact, known intervals, with position tracked from commanded steps rather than an encoder. This assumes that the motors do not miss steps. 

Two stepper motors are used in the gantry plotter, one for x and one for the y-axis direction. The type of stepper motor chosen was the 42BYGH40-01 Nema 17 bipolar stepper motor, with a resolution of 200 full steps per revolution and a current rating of 1.68A per phase~\cite{stepperds}. The following calculations were done to determine if the Nema 17 stepper motor would be strong enough for the use case of the gantry plotter, and if a singular motor could control both y-axis belt drives safely.

Dziedzic et al. examined samples written at 10--60\,mm/s~\cite{dziedzic2026}. The original sizing check uses the upper value, rounded from 3600 to 4000\,mm/min for headroom; this is a design benchmark, not the selected operating feed. Current generated jobs use 900\,mm/min for drawing and 2400\,mm/min for travel.

For the installed 20-tooth GT2 pulleys with 2\,mm belt pitch, the transmission geometry gives~\cite{acarnley2002}
\begin{equation}
\ell_{\mathrm{rev}}=N_t p=20(2)=40\ \mathrm{mm/rev},\qquad
r_p=\frac{N_t p}{2\pi}=6.37\ \mathrm{mm}.
\label{eq:des-pulley}
\end{equation}
\paragraph{Speed-to-torque margin check.}
The design linear feed is converted to a motor speed and compared with the torque required to move the assembly. The earlier curve readings are retained as a conditional calculation: their applicability to the installed motor and driver settings has not yet been verified.

\begin{equation}
n = \frac{v}{\ell_{\mathrm{rev}}} = \frac{4000\ \mathrm{mm/min}}{40\ \mathrm{mm/rev}} = 100\ \mathrm{RPM}
\label{eq:gan-rpm}
\end{equation}

where $n$ is the motor speed in revolutions per minute, $v=4000$\,mm/min is the sizing feed (60\,mm/s rounded up for margin), and $\ell_{\mathrm{rev}}=40$\,mm is the belt travel per revolution already established in Equation~\ref{eq:des-pulley} from the GT2 pulley's pitch ($p=2$\,mm, $N_t=20$ teeth).

\begin{equation}
T_{\mathrm{po}}(n) = T_1 + \frac{n-n_1}{n_2-n_1}\bigl(T_2-T_1\bigr) = 7 + \frac{100-60}{120-60}(21-7) \approx 16.3\ \mathrm{N\,cm}
\label{eq:gan-pullout}
\end{equation}

where $T_{\mathrm{po}}(n)$ is the conditional pull-out torque estimate at speed $n$, found by linear interpolation between the original two curve readings: $(n_1,T_1)=(60\ \mathrm{RPM},\,7\ \mathrm{N\,cm})$ and $(n_2,T_2)=(120\ \mathrm{RPM},\,21\ \mathrm{N\,cm})$. 

\begin{equation}
T_{\mathrm{req}} = r\,F_f = r\,\mu\,m\,g = 0.006366 \times 0.004 \times 0.950 \times 9.81 \approx 0.0237\ \mathrm{N\,cm}
\label{eq:des-torque}
\end{equation}

where $T_{\mathrm{req}}$ is the torque needed at constant velocity, $r=6.37$\,mm is the pulley pitch radius from Equation~\ref{eq:des-pulley}, $\mu=0.004$ is the original assumed bearing-friction coefficient, $m=0.950$\,kg is the complete moving assembly and $g=9.81$\,m/s$^2$. The bearings support the weight, so the motor overcomes horizontal resistance rather than lifting the assembly. The relations $F_f=\mu mg$ and $T=rF_f$ give this steady-speed estimate~\cite{acarnley2002}.

\begin{equation}
\frac{T_{\mathrm{po}}}{T_{\mathrm{req}}} = \frac{16.33}{0.02373} \approx 688
\label{eq:gan-margin}
\end{equation}

Under the assumed friction and original curve values, the estimated torque ratio is about 688. This retains the sizing argument for one motor driving both sides through a shared shaft. It is not a measured operating margin and does not include acceleration, preload or misalignment.

\begin{figure}[!htbp]
    \centering
    \begin{minipage}[c]{0.48\linewidth}
        \centering
        \includegraphics[width=\linewidth,height=0.2\textheight,keepaspectratio]{Figures/nema-17-stepper-motor.png}
    \end{minipage}\hfill
    \begin{minipage}[c]{0.48\linewidth}
        \centering
        \includegraphics[width=\linewidth,height=0.2\textheight,keepaspectratio]{Figures/a4988-stepper-driver.png}
    \end{minipage}
    \caption{NEMA 17 stepper motor (left) and A4988 stepper motor driver (right).}
    \label{fig:stepper-motor-and-driver}
\end{figure}

To power the two stepper motors, two A4988 stepper motor drivers are used.

\subsubsection{Step resolution, microstepping and calibration}
\label{sec:gan-res}
\label{sec:gan-timing}

Three properties of a stepper drive fix the resolution of this machine. The motors used here take 200 \emph{full steps} per revolution, so one full step is $1.8^\circ$; because position is simply the number of steps commanded, no encoder reports where the carriage actually is, which is what \emph{open-loop} control means, and the design simply trusts the step count. That trust is the assumption on which requirement R5 rests, and it fails silently if a step is ever missed. The driver subdivides each full step by \emph{microstepping}: rather than switching the two windings fully on or off, it holds them at intermediate current ratios, placing the rotor at intermediate angles. At the 1/16 setting used here, one full step becomes 16 microsteps, so a revolution is 3200 microsteps; microstepping buys resolution and smoother motion but not proportionally more accuracy, because the intermediate positions are held by the ratio of two winding currents and are therefore more sensitive to load torque than a full step is.

Equation~\ref{eq:des-stepres} gives the step resolution from 200 full steps~\cite{stepperds}, 16 microsteps~\cite[p.~6]{a4988ds} and $\ell_{\mathrm{rev}}=40$\,mm,
\begin{equation}
\kappa=\frac{200\cdot16}{\ell_{\mathrm{rev}}}=\frac{3200}{40}=80\ \mathrm{steps/mm},\qquad\Delta=\frac1\kappa=12.5\ \mu\mathrm{m},
\label{eq:des-stepres}
\end{equation}
which calibration replaces by each axis's measured value. Targets are \emph{absolute} step counts, $\mathrm{round}\bigl((x+5)\kappa\bigr)$ with a 5\,mm offset, differenced against the tracked position, so rounding cannot accumulate. Quantisation error is at most $\Delta/2=6.25\ \mu$m per axis, $8.8\ \mu$m combined, 57 times finer than the 0.5\,mm of R5 --- but that is the \emph{digital} resolution only, and belt stretch, backlash, missed steps, switch repeatability and pen compliance were not measured.
\paragraph{Coil current and its derating.}
The datasheet relation~\cite[p.~9]{a4988ds} between the trimmer reference, the carrier's two sense resistors and the regulated current gives
\begin{equation}
I_{\mathrm{coil}}=\frac{V_{\mathrm{ref}}}{8\,R_S}\quad\Longrightarrow\quad V_{\mathrm{ref}}=8\,R_S\,I_{\mathrm{coil}}=8\cdot0.1\cdot1.68=1.34\ \mathrm{V}
\label{eq:des-current}
\end{equation}
for the 1.68\,A motor rating~\cite{stepperds}, conditional on the reported R10 sense-resistor marking meaning $0.1\,\Omega$. This gives a rating-based ceiling, rather than the measured adjustment. For comparison, reducing the phase-current limit to 1\,A would give $V_{\mathrm{ref}}=8(0.1)(1)=0.8$\,V. Carrier cooling can require a lower current, but the actual fitted $V_{\mathrm{ref}}$ and thermal margin were not recorded.

\subsubsection{Pen up/down mechanism}
\label{sec:des-pen}\label{sec:gan-circuits}\label{sec:des-power}
The pen-lift mechanism itself is a design choice justified against the alternative of a third stepper or a servo motor. A servo motor expects a continuous pulse-width-modulated position command and typically needs its own timer peripheral or a precisely timed software loop running concurrently with the gantry's own step timing, competing for the same processor cycles that generate STEP pulses. An N20 geared d.c.\ motor run in one direction only, with a cam that toggles the pen up or down every $90^\circ$ of rotation, needs only an on/off drive pulse of roughly known duration and a microswitch to confirm the result; it is cheap, widely available, has high torque at low speed through its internal gearbox for its size, and has a smaller profile than a comparable servo, at the cost of giving no absolute position, only a toggle.

\begin{figure}[!htbp]
    \centering
    \includegraphics[width=0.75\linewidth]{Figures/pen-up-down-mechanism.png}
    \caption{Cam-operated pen-lift mechanism in the down (left) and up (right) positions.}
    \label{fig:pen-up-down-mechanism}\label{fig:des-penmechanism}
\end{figure}

The implemented lift, planned in the proposal as a stepper, gives only pulses, so the pen state must be tracked in software, and the pen switch is feedback only (high = up). The pen-drive routine of Algorithm~\ref{alg:gan-pen} must therefore test the switch \emph{before} energising the motor, since the cam turns one way only and an unnecessary pulse would toggle the pen to the wrong state rather than leave it alone. The controller runs the motor until the requested switch state is confirmed or a timeout occurs, with a 40\,ms overrun after pen-up confirmation. The generator's separate settling dwell is 350\,ms; this is not a measured fixed motor-rotation time.

\paragraph{Pen motor driver and base-resistor selection.}
The GA12-N20 geared motor was selected to turn the pen cam with a compact, low-speed geared output. The listed 3\,V, 150\,rpm variant has a specified voltage range of 1.5--5\,V, a no-load current of 10\,mA, a rated current of 30\,mA, and a stall current of 230\,mA; these are supplier figures for that variant, rather than measured currents under the installed cam load~\cite{ga12n20ds}. The supplied circuit uses the Odroid's 3.3\,V GPIO signal to drive the base of a BC337 NPN transistor. The transistor is a low-side switch: motor current flows from the 3.3\,V DC motor rail, through the motor and collector-emitter path to the shared ground, while the GPIO supplies only base current~\cite{odroidn2ds}. The source symbol in the supplied diagram should therefore be read as a DC supply, not an AC source.

Assuming a base-emitter voltage of approximately 0.7\,V, the 1\,k$\Omega$ base resistor limits the nominal GPIO current to
\begin{equation}
I_B = \frac{V_{\mathrm{GPIO}}-V_{BE}}{R_B} = \frac{3.3-0.7}{1000} \approx 2.6\,\mathrm{mA}.
\label{eq:des-basecurrent}
\end{equation}
This keeps the GPIO current far below the motor current, but the GPIO must still be rated to source at least 2.6\,mA. The BC337 datasheet's saturation test uses 50\,mA of base current for 500\,mA of collector current, a conservative forced-current-gain ratio of 10~\cite{bc337ds}. Using that ratio as a design check, 2.6\,mA of base current supports about 26\,mA of saturated collector current. That is slightly below the listed 30\,mA rated motor current. To provide the same conservative saturation drive at 30\,mA would require $R_B\le(3.3-0.7)/(0.030/10)\approx870\,\Omega$; at the listed 230\,mA stall current it would require $R_B\le(3.3-0.7)/(0.230/10)\approx110\,\Omega$. The latter would draw about 23\,mA from the GPIO, so the 1\,k$\Omega$ circuit is not guaranteed to saturate the transistor at rated or stall load. The 1\,k$\Omega$ resistor is the prototype's GPIO-current-limiting choice. A suitable logic-level MOSFET or motor driver would avoid relying on this base-current margin if the installed motor approaches those loads.

The supplied circuit image does not show a flyback diode. A diode should be connected across the motor, reverse-biased during normal operation, with its cathode toward the positive motor rail and anode at the transistor collector, to provide a path for the inductive current when the transistor switches off~\cite{rohmlowside2026}.

\begin{figure}[!htbp]
    \centering
    \includegraphics[width=0.68\linewidth]{Figures/pen-clicker-bc337-circuit.png}
    \caption{Supplied low-side BC337 switch circuit for the GA12-N20 pen-lift motor.}
    \label{fig:pen-clicker-bc337-circuit}\label{fig:des-pencircuit}
\end{figure}

\begin{algorithm}[!htbp]
\caption{Pen toggle}
\label{alg:gan-pen}
\begin{algorithmic}[1]
\REQUIRE desired state $\mathrm{wantUp}$
\IF{any end-stop is currently triggered}
    \STATE refuse the request
\ELSIF{the pen switch already reads $\mathrm{wantUp}$}
    \STATE \textbf{return} \COMMENT{the motor must \emph{not} run; the cam turns one way only}
\ELSE
    \STATE energise the pen motor
    \WHILE{before the deadline and the pen switch does not yet read $\mathrm{wantUp}$}
        \STATE wait
    \ENDWHILE
    \IF{the switch now reads up}
        \STATE wait 40\,ms more to seat the cam
    \ENDIF
    \STATE de-energise the pen motor
\ENDIF
\end{algorithmic}
\end{algorithm}

During tuning the clicker sometimes over-rotated and the pen did not push down fully before the switch read as triggered, requiring the switch seating and timing to be adjusted experimentally. Differing pen pressure on a book or an exam pad also introduced inconsistent ink, so a rigid clipboard with a single page clamped flat was adopted. The 0.35\,s settling dwell contributes to the predicted writing time, but is separate from switch-confirmed actuation. Missed steps cannot be detected from the step counter, so physical accuracy must be assessed from the drawn output rather than nominal resolution alone.

\subsubsection{End stop switches and calibration}
\label{sec:des-calibration}
All four axis end-stop switches are checked before every major-axis step and every G-code line. A potentially triggered switch is accepted only after three consecutive high reads 6\,ms apart (Algorithm~\ref{alg:gan-endstop}). The pen switch is read separately by the pen-control routine. The confirmation interval is
\begin{equation}
t_{\mathrm{confirm}}=(3-1)(6\,\mathrm{ms})=12\,\mathrm{ms}.
\label{eq:gan-latency}
\end{equation}
The stepping loop does not continue pulsing during this check. The interval is therefore not a measured trip-to-stop latency and cannot simply be multiplied by the requested feed to infer stopping distance.

The DML31P was selected as a compact, momentary lever microswitch for detecting gantry end-stop contact; the product listing identifies it as an SPDT switch rated for 1\,A at 125\,V AC~\cite{dml31p}. In this circuit it is used only as a low-current logic contact, not to switch a motor. The diagram uses the common and normally-open contacts so that actuating the switch connects the Odroid N2+'s 3.3\,V logic rail to the GPIO input; the source symbol represents DC. The 10\,k$\Omega$ resistor is a pull-down: it holds the input low while the switch is open, preventing a floating input, and draws a small current when the switch closes. The N2+ GPIO uses 3.3\,V logic~\cite{odroidn2ds}.

With the switch closed, the pull-down current and resistor dissipation are
\begin{equation}
I_{\mathrm{closed}}=\frac{V_{\mathrm{rail}}}{R_{\mathrm{PD}}}=\frac{3.3\,\mathrm{V}}{10\,\mathrm{k}\Omega}=0.33\,\mathrm{mA},\qquad
P_{\mathrm{PD}}=\frac{V_{\mathrm{rail}}^2}{R_{\mathrm{PD}}}=\frac{(3.3\,\mathrm{V})^2}{10\,\mathrm{k}\Omega}\approx1.1\,\mathrm{mW}.
\label{eq:endswitch-pulldown}
\end{equation}
The 0.33\,mA flows from the 3.3\,V rail through the switch and resistor to ground; only the GPIO input leakage flows into the pin. When the switch opens, the resistor pulls the input to ground; when it closes, the input rises to approximately 3.3\,V. This provides defined logic levels while keeping the closed-contact current low.
The GPIO must be configured as an input before this circuit is connected and kept in input mode during operation. If the pin were configured as an output-low, closing the switch would connect the 3.3\,V rail directly to the GPIO output driver; the $10\,k\Omega$ pull-down does not limit that fault current.

\begin{figure}[!htbp]
    \centering
    \includegraphics[width=0.75\linewidth]{Figures/endstop-switch-pulldown-circuit.png}
    \caption{DML31P end-stop switch with a 10\,k$\Omega$ pull-down connected to an Odroid N2+ GPIO input.}
    \label{fig:endstop-switch-pulldown}\label{fig:des-switch}
\end{figure}

\begin{algorithm}[!htbp]
\caption{Debounced end-stop read, called before every major-axis pulse}
\label{alg:gan-endstop}
\begin{algorithmic}[1]
\REQUIRE set of end-stops to ignore (already triggered on this axis)
\FOR{each end-stop not in the ignore set}
    \IF{it does not read closed}
        \STATE \textbf{continue} to the next end-stop
    \ENDIF
    \STATE re-read it twice more, 6\,ms apart \COMMENT{three reads in total, debounce of Equation~\ref{eq:gan-latency}}
    \IF{it read closed on all three reads}
        \STATE \textbf{return} that end-stop as triggered
    \ENDIF
\ENDFOR
\STATE \textbf{return} none triggered
\end{algorithmic}
\end{algorithm}

On a trip the program latches a stop, sets both STEP pins low, retracts off every triggered switch and aborts the job; targets are clamped to the calibrated area and the latch clears only by explicit command, refused while a switch is still triggered. Protection is \emph{software only}: no hardware stop that cuts motor power is documented, the switches are polled only during a job, the web stop request can interrupt a running job, but it is not an independent hardware power cut.

\label{sec:gan-calib}
Equation~\ref{eq:des-stepres} rests on assumed transmission values and the position after power-up is unknown, so each session begins with a direction-agnostic calibration, written so that it never needs to know which direction leads to which switch. It drives one way at 1.6\,ms per step until \emph{either} switch of that axis triggers, reverses, and counts the steps $N$ to the \emph{other}. With the hand-measured travel $L$, the resolution, the initial step counter and the position in millimetres are
\begin{equation}
\kappa_{\mathrm{axis}}=\frac{N}{L},\qquad n_0=N\ \text{if the second switch was MAX else }0,\qquad x_{\mathrm{mm}}=\frac{n}{\kappa}-5,
\label{eq:des-calibration}
\end{equation}
so the zero of Equation~\ref{eq:des-calibration} lies 5\,mm inside the minimum switch, the usable area is $190.6\times251.3$\,mm, and each axis retracts 5\,mm off its last switch. Using the nominal 80 microsteps/mm, a 1.6\,ms homing-step period corresponds to 7.8125\,mm/s. The counted switch-to-switch traverses alone would then need $(200.6+261.3)/7.8125\approx59.1$\,s; this excludes the initial searches, retractions and parking, and is not a measured calibration duration. The result is deliberately not stored between sessions, since a bump between sessions would silently invalidate it, and a failed calibration leaves movement unavailable until calibration succeeds. Algorithm~\ref{alg:gan-calibrate} states the procedure.

\begin{algorithm}[!htbp]
\caption{Direction-agnostic calibration of one axis}
\label{alg:gan-calibrate}
\begin{algorithmic}[1]
\STATE drive one way at 1.6\,ms per step until \emph{either} switch of this axis triggers
\STATE reverse; count steps $N$ until the \emph{other} switch triggers
\STATE $\kappa_{\mathrm{axis}}\leftarrow N/L$ \COMMENT{$L$ the hand-measured travel, Equation~\ref{eq:des-calibration}}
\IF{the second switch to trigger was the maximum-side one}
    \STATE $n_0\leftarrow N$
\ELSE
    \STATE $n_0\leftarrow0$
\ENDIF
\STATE retract 5\,mm off the switch that is still closed
\IF{an already-triggered switch, an unexpected switch on the other axis, or 200\,000 steps pass with no trigger}
    \STATE abort calibration; refuse all jobs until it is re-run
\ENDIF
\end{algorithmic}
\end{algorithm}

The hand measurement of travel also sets the scale of every subsequent movement. An error $\delta L$ scales every commanded position by $x\,\delta L/L$, so to keep that scale error below half of the 0.5\,mm allowed by R5 over the 190.6\,mm of usable X travel,
\begin{equation}
\delta L\le\frac{0.25\ \mathrm{mm}}{190.6\ \mathrm{mm}}\,L=\frac{0.25}{190.6}\cdot200.6\ \mathrm{mm}=0.26\ \mathrm{mm}\quad(0.13\ \%),
\label{eq:gan-budget}
\end{equation}
and switch-trigger repeatability adds directly to the zero error on top of that. This is an allocation of half the positional-error allowance to scale error, not a measured uncertainty budget. Switch repeatability and mechanical compliance use part of the remaining allowance~\cite[pp.~18--19]{jcgm2008}.

A mechanical switch does not change state cleanly: its contacts strike, separate and strike again over a few milliseconds, so a single read taken during that interval may return either level. Treating one high read as a trip would stop the machine at random, and treating the first read as final would miss a genuine trip. The standard remedy is to \emph{debounce} the input, requiring the same level on several reads spaced far enough apart to outlast the bouncing.
\subsubsection{Gcode generation and motion mapping}
\label{sec:des-drive}
G-code stores movement and pen-state instructions in a text file. The hardware controller reads each command and its parameters, then converts the target coordinates to motor pulses. The instructions used by this plotter are explained below before the square-and-triangle commissioning example.

\paragraph{Interpolating a straight line.}
Stepping the two axes independently would draw an L-shaped path, so the program interpolates with Bresenham's integer line algorithm~\cite[p.~28]{bresenham1965}: an error term $e$ is carried along the major axis of $a_x=|\Delta x|$ and $a_y=|\Delta y|$ steps, the major axis is pulsed every iteration and the minor axis whenever $e$ would go negative, and the end-stops are read before every major step. Keeping $e$ in $[0,a_x)$ gives the minor-axis count after $k$ major steps as
\begin{equation}
m_k=\Bigl\lceil\frac{k\,a_y-\lfloor a_x/2\rfloor}{a_x}\Bigr\rceil,
\label{eq:gan-bres}
\end{equation}
which is within half a microstep of the ideal $k\,a_y/a_x$, at most 6.25\,\textmu m at the nominal 80 microsteps/mm, using integer arithmetic only. Equation~\ref{eq:gan-bres} is the reason no floating-point division appears in the inner loop: the deviation from the true line is bounded by construction rather than by rounding behaviour. Algorithm~\ref{alg:gan-step} states the loop.

\begin{algorithm}[!htbp]
\caption{Coordinated line movement with an integer error accumulator}
\label{alg:gan-step}\label{alg:des-stepping}
\begin{algorithmic}[1]
\REQUIRE signed target changes in microsteps
\STATE set axis directions; choose major count $N_M$ and minor count $N_m$
\STATE $e\leftarrow\lfloor N_M/2\rfloor$
\FOR{each of the $N_M$ major-axis steps}
    \IF{abort requested or an end-stop is confirmed}
        \STATE \textbf{return} interrupted
    \ENDIF
    \STATE pulse major axis; update its count; $e\leftarrow e-N_m$
    \IF{$e<0$}
        \STATE pulse minor axis; update its count; $e\leftarrow e+N_M$
    \ENDIF
\ENDFOR
\end{algorithmic}
\end{algorithm}

Figure~\ref{fig:vt-stepping} illustrates the commanded microsteps for a separate generated move of $(+6,-28)$ microsteps. Its maximum deviation from the ideal line is 0.5 microsteps, or $0.5/80=0.00625$\,mm at the nominal scale. This checks command-space interpolation; it does not measure the pen's physical path.
\begin{figure}[!htbp]\centering
\includegraphics[width=0.96\linewidth]{Figures/VisualTests/T52_bresenham_staircase.pdf}
\caption{Integer axis stepping and command-space deviation for a generated move.}
\label{fig:vt-stepping}
\end{figure}

Equation~\ref{eq:des-pulsetime} gives the delay of every pulse for a move of length $d$ at feed $v_f$ over $N_{\mathrm{maj}}=\max(|n_x|,|n_y|,1)$ major-axis steps,
\begin{equation}
\tau=\max\Bigl(\frac{d/v_f}{N_{\mathrm{maj}}},\ 2\,\mu\mathrm s\Bigr),\qquad v_f=\max(F/60,\,0.1),
\label{eq:des-pulsetime}
\end{equation}
with $F$ in mm/min: 900 gives $\tau=833\ \mu$s (1200 steps/s), 2400 gives 312.5\,\textmu s (3200 steps/s, 60\,rpm), homing 625 steps/s or 7.8\,mm/s. Because \emph{every} pulse of both axes sleeps for $\tau$, a segment takes $(n_{\mathrm{maj}}+n_{\mathrm{min}})\tau$ and Equation~\ref{eq:gan-eff} gives the effective feed,
\begin{equation}
v_{\mathrm{eff}}=v_f\,\frac{n_{\mathrm{maj}}}{n_{\mathrm{maj}}+n_{\mathrm{min}}}\in\Bigl[\frac{v_f}{2},v_f\Bigr],
\label{eq:gan-eff}
\end{equation}
unity axis-parallel, $\tfrac12$ at $45^\circ$. There are no controlled acceleration ramps; the steady-speed torque estimate alone does not establish a safe starting rate~\cite{acarnley2002}.


Figure~\ref{fig:gcode-square-triangle-comparison} compares the generated path preview with the corresponding physical plot. The test file commands a 100\,mm square and an 80\,mm-base triangle, providing both horizontal and diagonal movements for visual comparison.

\begin{figure}[!htbp]
    \centering
    \begin{minipage}[c]{0.48\linewidth}
        \centering
        \includegraphics[width=\linewidth,height=0.30\textheight,keepaspectratio]{Figures/square-triangle-gcode-preview.png}\\
        (a) G-code path preview
    \end{minipage}\hfill
    \begin{minipage}[c]{0.48\linewidth}
        \centering
        \includegraphics[width=\linewidth,height=0.30\textheight,keepaspectratio]{Figures/square-triangle-reproduced.jpg}\\
        (b) Reproduced drawing
    \end{minipage}
    \caption{Comparison of the G-code preview and the physical square-and-triangle plot.}
    \label{fig:gcode-square-triangle-comparison}\label{fig:des-shapetest}
\end{figure}

The following extract shows the commands used to draw the square; the full file then continues with the triangle before its final \texttt{M2}. Coordinates are in millimetres and use absolute positioning; the pen is lowered before the drawing moves and raised after the square is complete.

\begin{verbatim}
G21
G90
G0 F2400
G0 X45.000 Y76.000
M3
G4 P0.350
G1 F900
G1 X145.000 Y76.000
G1 X145.000 Y176.000
G1 X45.000 Y176.000
G1 X45.000 Y76.000
M5
G4 P0.350
\end{verbatim}

Table~\ref{tab:gcode-command-meanings} describes the commands and parameters used in this file. Travel commands do not themselves raise the pen, and the travel feed is restated after drawing because \texttt{F} remains active until changed. The meanings of \texttt{M3} and \texttt{M5} are specific to this plotter firmware, where they operate the pen-lift mechanism rather than a spindle.

\begin{table}[!htbp]
    \centering
    \begin{tabularx}{\linewidth}{|l|X|}
        \hline
        \textbf{Command or parameter} & \textbf{Meaning in this plotter} \\
        \hline
        \texttt{G21} & Select millimetres as the units for coordinates and feed rates. \\
        \hline
        \texttt{G90} & Use absolute positioning, so each X/Y coordinate is measured from the machine origin. \\
        \hline
        \texttt{G0} & Move rapidly without drawing, used for positioning and travel with the pen raised. \\
        \hline
        \texttt{G1} & Make a straight, feed-controlled move; with the pen down, this draws a line. \\
        \hline
        \texttt{G4 P} & Dwell for the duration specified by P in seconds; used here to let the pen-lift motor finish moving. \\
        \hline
        \texttt{M3} & Lower the pen using the plotter's one-way pen-lift motor. \\
        \hline
        \texttt{M5} & Raise the pen using the plotter's one-way pen-lift motor. \\
        \hline
        \texttt{M2} & End the G-code program. \\
        \hline
        \texttt{X}, \texttt{Y} & Target coordinates in millimetres when used with the active positioning mode. \\
        \hline
        \texttt{F} & Feed rate for movement, in millimetres per minute with \texttt{G21} active. \\
        \hline
        \texttt{P} & Dwell duration in seconds when supplied to \texttt{G4}. \\
        \hline
    \end{tabularx}
    \caption{G-code commands and parameters used in the square-and-triangle test file.}
    \label{tab:gcode-command-meanings}
\end{table}


\FloatBarrier
